\documentclass[conference]{IEEEtran}
\IEEEoverridecommandlockouts
% The preceding line is only needed to identify funding in the first footnote. If that is unneeded, please comment it out.
\usepackage{cite}
\usepackage{amsmath,amssymb,amsfonts}
\usepackage{algorithmic}
\usepackage{graphicx}
\usepackage{textcomp}
\usepackage{xcolor}

%------Newly added packages 25-06-2026
\usepackage{booktabs}
\usepackage{multirow}
\usepackage{subfigure}
\usepackage{caption}
\usepackage{algorithm,algorithmic}
\usepackage{gensymb}
\usepackage{bm}
\usepackage{lipsum}
%--------------------------------

\def\BibTeX{{\rm B\kern-.05em{\sc i\kern-.025em b}\kern-.08em
    T\kern-.1667em\lower.7ex\hbox{E}\kern-.125emX}}
\begin{document}

\title{Deep Learning-Based Security Framework for Internet of Things: A Survey}

\author{\IEEEauthorblockN{Ram Kapadia}
\IEEEauthorblockA{\textit{Student ID: 2548530}\\
\textit{M.Sc. Artificial Intelligence \& Machine Learning}\\
\textit{CHRIST (Deemed to be University)}\\
\textit{Course: IoT Analytics} \\
\textit{Faculty: Dr. Jobin Francis}
}
}

\maketitle

\begin{abstract}
The rapid proliferation of the Internet of Things (IoT) has transformed various domains, from smart healthcare to industrial automation, by interconnecting billions of resource-constrained devices. However, this expansive growth has broadened the attack surface, rendering IoT networks susceptible to sophisticated cyber threats such as zero-day attacks, distributed denial-of-service (DDoS), and data breaches. Traditional security mechanisms, relying primarily on static rules and signature-based detection, are often inadequate due to the heterogeneous and dynamic nature of IoT environments. Consequently, Deep Learning (DL) has emerged as a promising paradigm to develop intelligent, adaptive, and scalable security frameworks capable of automatically extracting complex features from massive volumes of network traffic. This survey provides a comprehensive review of recent DL-based security solutions for IoT, evaluating architectures like Convolutional Neural Networks (CNNs), Recurrent Neural Networks (RNNs), Long Short-Term Memory (LSTM) networks, Autoencoders, Transformers, and Federated Learning. By systematically analyzing state-of-the-art literature published between 2021 and 2025 across Q1 journals, we examine the methodologies, datasets, advantages, and limitations of these approaches. Furthermore, this paper identifies critical research gaps concerning model explainability, privacy, resource constraints, and real-time deployment, ultimately proposing future research directions to advance the field of artificial intelligence-driven IoT security.
\end{abstract}

\begin{IEEEkeywords}
Internet of Things, Deep Learning, Artificial Intelligence, IoT Security, Intrusion Detection, Cybersecurity, Machine Learning.
\end{IEEEkeywords}

\section{Introduction}
The Internet of Things (IoT) refers to the vast, interconnected network of physical devices equipped with sensors, software, and connectivity, enabling them to collect and exchange data continuously. This technological paradigm has fundamentally reshaped numerous sectors, optimizing operations in smart cities, industrial control systems (ICS), and personalized healthcare. However, the ubiquitous deployment of IoT devices presents unprecedented security challenges. Because IoT nodes are often resource-constrained—possessing limited computational power, memory, and energy—they are frequently deployed with weak or default security configurations. This makes them prime targets for malicious actors seeking to launch large-scale cyberattacks, such as Distributed Denial of Service (DDoS) campaigns, botnet infections, and ransomware attacks.

Traditional intrusion detection systems (IDS) and firewalls typically rely on signature-based methodologies, which match network traffic against a database of known threat profiles. While effective against recognized attacks, these systems fail to detect novel, zero-day vulnerabilities. Furthermore, conventional security techniques struggle to handle the dynamic nature, heterogeneous data streams, and massive scale typical of modern IoT networks.

In response to these challenges, Deep Learning (DL), a sophisticated subset of Artificial Intelligence (AI), offers a robust alternative. DL models, such as Convolutional Neural Networks (CNNs) and Recurrent Neural Networks (RNNs), possess the intrinsic ability to autonomously learn hierarchical representations of data directly from raw network packets. This allows for the highly accurate detection of subtle anomalies and complex intrusion patterns without explicit, labor-intensive feature engineering. 

This survey systematically explores the intersection of deep learning and IoT security. It evaluates the efficacy of various DL architectures, identifies current research gaps, and outlines the trajectory for future intelligent security frameworks designed specifically for the rigorous demands of IoT ecosystems.

\section{Problem Statement}
Despite substantial advancements in cybersecurity, current IoT security systems continually fail to mitigate evolving cyber threats effectively. The primary issue stems from the heterogeneity and resource constraints of IoT devices, which prevent the deployment of heavy, traditional cryptographic protocols and conventional security agents. Furthermore, as attack vectors become increasingly sophisticated—leveraging automation and AI to launch zero-day exploits and polymorphic malware—static, rule-based systems are rendered obsolete. There is an urgent, unmet need for lightweight, intelligent, and adaptive security solutions that can autonomously detect, classify, and mitigate threats in real-time, directly at the network edge, without compromising device performance or data privacy.

\section{Objectives}
The primary objectives of this survey are to:
\begin{itemize}
    \item \textbf{Study existing IoT security threats:} Analyze the evolving landscape of cyberattacks targeting IoT infrastructures.
    \item \textbf{Review deep learning approaches:} Comprehensively evaluate how different DL architectures are applied to intrusion and anomaly detection.
    \item \textbf{Compare distinct methodologies:} Provide a critical comparison of CNN, RNN, LSTM, Autoencoder, Transformer, and Federated Learning methods.
    \item \textbf{Identify research gaps:} Pinpoint the limitations of current state-of-the-art models, focusing on explainability, privacy, and computational overhead.
    \item \textbf{Suggest future research directions:} Outline actionable pathways for developing scalable, robust, and lightweight AI-driven IoT security frameworks.
\end{itemize}

\section{Survey Methodology}
To ensure a rigorous and comprehensive review, this survey adopted a systematic methodology consisting of the following sequential phases:
\begin{enumerate}
    \item \textbf{Literature Collection:} We conducted extensive searches across major academic databases, including IEEE Xplore, ACM Digital Library, Elsevier ScienceDirect, and SpringerLink.
    \item \textbf{Paper Selection:} We filtered the collected literature to include peer-reviewed articles published in Q1 journals between 2021 and 2025.
    \item \textbf{Categorization:} The selected papers were categorized based on the underlying DL architecture.
    \item \textbf{Comparative Analysis:} We critically evaluated each approach against key metrics, including the utilized dataset, detection accuracy, advantages, and limitations.
    \item \textbf{Research Gap Identification:} By synthesizing the limitations, we identified prevailing challenges in the field.
    \item \textbf{Future Directions:} Based on the identified gaps, we proposed strategic avenues for future research.
\end{enumerate}

\section{Literature Review}
The integration of DL in IoT security has seen substantial research interest. This section reviews notable studies published between 2021 and 2025.

Liu \textit{et al.} \cite{liu2021novel} proposed a novel IDS utilizing Convolutional Neural Networks (CNNs) for spatial feature extraction in IoT networks. Evaluated on the NSL-KDD dataset, the model achieved exceptional accuracy in detecting DDoS attacks. The primary advantage of this approach was its rapid inference time; however, the model struggled to capture long-term temporal dependencies, limiting its effectiveness against slow-rate stealth attacks.

Zhao \textit{et al.} \cite{zhao2021} addressed temporal sequence modeling by introducing an RNN framework. Tested on the Bot-IoT dataset, the model successfully identified sequentially structured botnet attacks. While providing excellent temporal awareness, it suffered from the vanishing gradient problem during prolonged traffic sequences and required significant computational overhead, complicating edge deployment.

To overcome the limitations of standard RNNs, Kim and Park \cite{kim2023} developed an anomaly detection system integrating Autoencoders with Temporal Convolutional Networks (TCN). Tested against the CICIDS2017 dataset, the unsupervised Autoencoder efficiently reconstructed normal traffic profiles, while the TCN flagged deviations. This hybrid approach excelled in identifying zero-day attacks without labeled data but exhibited higher false positive rates in highly dynamic environments.

Zhang \textit{et al.} \cite{zhang2022} introduced a lightweight Transformer-based IDS designed for resource-constrained IoT nodes. Utilizing the ToN\_IoT dataset, the self-attention mechanism allowed the model to focus on critical packet headers. The framework demonstrated superior accuracy and lower latency compared to traditional LSTMs. Nonetheless, the high memory footprint during the training phase remained a significant constraint for complete on-device execution.

Alqahtani \textit{et al.} \cite{alqahtani2022} and Ali and Khan \cite{ali2024} explored privacy-preserving threat detection through Federated Learning (FL). By allowing multiple IoT gateways to train a shared DL model collaboratively without exchanging raw data, they successfully detected distributed attacks. The approach significantly enhanced data privacy; however, high communication overhead and susceptibility to data poisoning attacks were identified as major limitations.

Ahanger \textit{et al.} \cite{ahanger2024} implemented a comprehensive DL-based technique emphasizing real-time feature extraction. Their evaluation on the BoT-IoT dataset demonstrated high precision, though the requirement for continuous model retraining poses a barrier to autonomous operation. Finally, Wang \textit{et al.} \cite{wang2025} proposed a self-supervised deep learning approach specifically targeting industrial IoT, highlighting robustness against zero-day attacks, though its complexity hinders deployment on low-power sensor nodes.

\section{Comparative Analysis}
Table \ref{tab:comparison} provides a comprehensive overview of the recent deep learning-based IoT security frameworks reviewed in this survey.

\begin{table*}[htbp]
\centering
\caption{Comparative Analysis of Recent Deep Learning-Based IoT Security Approaches}
\label{tab:comparison}
\begin{tabular}{@{}llp{2cm}lp{1.2cm}p{3.5cm}p{3.5cm}@{}}
\toprule
\textbf{Author} & \textbf{Year} & \textbf{Technique} & \textbf{Dataset} & \textbf{Accuracy} & \textbf{Advantages} & \textbf{Limitations} \\ \midrule
Liu \textit{et al.} \cite{liu2021novel} & 2021 & CNN & NSL-KDD & 97.5\% & High accuracy; rapid inference time for spatial features. & Fails to capture long-term temporal dependencies. \\
Zhao \textit{et al.} \cite{zhao2021} & 2021 & RNN & Bot-IoT & 96.8\% & Excellent temporal sequence modeling capabilities. & Vanishing gradient problem; high computational overhead. \\
Zhang \textit{et al.} \cite{zhang2022} & 2022 & Transformer & ToN\_IoT & 98.2\% & Low latency; effectively isolates critical packet headers. & High memory footprint during the training phase. \\
Alqahtani \textit{et al.} \cite{alqahtani2022} & 2022 & Federated Learning & IoT-23 & 95.4\% & Preserves data privacy; enables distributed threat detection. & High communication overhead; vulnerable to data poisoning. \\
Kim \textit{et al.} \cite{kim2023} & 2023 & Autoencoder + TCN & CICIDS2017 & 94.1\% & Unsupervised zero-day detection without labeled data. & High false positive rates in highly dynamic environments. \\
Ahanger \textit{et al.} \cite{ahanger2024} & 2024 & Hybrid DNN & BoT-IoT & 98.9\% & Comprehensive real-time feature extraction; high precision. & Requires frequent model retraining to maintain efficacy. \\
Wang \textit{et al.} \cite{wang2025} & 2025 & Self-Supervised DL & Industrial IoT & 97.8\% & Exceptional robustness against novel zero-day attacks. & Architectural complexity hinders low-power deployment. \\ \bottomrule
\end{tabular}
\end{table*}

\section{Research Gap}
Despite the promising results demonstrated by DL models, several critical research gaps remain unresolved:
\begin{itemize}
    \item \textbf{Explainability and Trust:} Most high-performing DL models act as "black boxes." The lack of Explainable AI (XAI) \cite{rahman2022} makes it difficult for security analysts to understand the rationale behind intrusion alerts, hindering trust and practical adoption.
    \item \textbf{Resource Constraints vs. Model Complexity:} State-of-the-art architectures (e.g., Transformers, deep CNNs) require substantial memory and processing power, which directly conflicts with the stringent limitations of edge IoT devices.
    \item \textbf{Privacy Concerns:} Centralized training mandates the transfer of massive amounts of sensitive network data to the cloud, posing severe privacy risks and increasing bandwidth consumption.
    \item \textbf{Federated Learning Challenges:} While FL mitigates data privacy issues, it introduces new challenges, including high communication overhead, synchronization latency, and vulnerability to adversarial data poisoning.
    \item \textbf{Lightweight Models:} Developing ultra-lightweight DL models that offer high detection accuracy without depleting device batteries remains an open challenge.
    \item \textbf{Zero-Day Attack Detection:} Many supervised models suffer from drastic performance degradation when encountering novel attack signatures not present in their training datasets.
    \item \textbf{Real-Time Deployment:} Achieving ultra-low latency inference for real-time threat mitigation remains a significant hurdle for complex DL pipelines deployed in dynamic, large-scale IoT networks.
\end{itemize}

\section{Expected Outcomes}
This survey comprehensively maps the current landscape of AI-based IoT security. By identifying prevailing trends, we highlight that the shift toward edge AI and decentralized learning is inevitable. However, the outcomes also underscore that raw accuracy is no longer the sole metric of success; deployability, explainability, and resource efficiency are equally paramount. The synthesis of current challenges will serve as a foundational guide for researchers, facilitating the development of the next generation of lightweight, privacy-preserving, and adaptive intrusion detection systems tailored for the IoT ecosystem.

\section{Conclusion}
The integration of Deep Learning into IoT security frameworks represents a critical advancement in defending against increasingly sophisticated cyber threats. This survey systematically reviewed recent literature from 2021 to 2025, evaluating various architectures ranging from CNNs and RNNs to Transformers and Federated Learning paradigms. While these models exhibit exceptional accuracy and the ability to detect complex anomalies, their practical deployment is hampered by issues related to computational overhead, lack of explainability, and privacy concerns. Future research must prioritize the development of Explainable AI (XAI) to foster trust, alongside the creation of ultra-lightweight models capable of executing directly on edge devices. Furthermore, enhancing the resilience of decentralized learning approaches against adversarial attacks will be crucial in realizing fully autonomous, secure, and scalable IoT environments.

\bibliographystyle{IEEEbib}
\bibliography{refs}

\end{document}
